\section{Preliminaries}
\label{sec:preliminaries}

This section reviews the two foundations of Kimodo-Policy. We first summarize Kimodo's hierarchical text-to-motion formulation and kinematic representation (Sec.~\ref{sec:prelim_kimodo}), and then introduce ControlNet and its parameter-efficient mechanism for injecting additional conditions into a frozen diffusion backbone (Sec.~\ref{sec:prelim_controlnet}).

\subsection{Kimodo}
\label{sec:prelim_kimodo}
\textbf{Architecture Overview.} Kimodo~\cite{kimodo} is a framework for controllable kinematic motion generation based on diffusion, with pretrained models for both the SOMA~\cite{saito2026somaunifyingparametrichuman} human skeleton and the Unitree G1 humanoid robot. Starting from Gaussian noise, it generates motion sequences conditioned on text prompts and optional kinematic constraints using a two-stage Transformer denoiser operating on pose tokens. At each diffusion step $t$, its input--output mapping is summarized as
\begin{equation}
  \widehat{\mathbf{X}}_0 = \mathcal{D}_{\theta}(\mathbf{X}_t,c,t)
  = [\widehat{\mathbf{R}}_0,\widehat{\mathbf{B}}_0].
  \label{eq:kimodo_two_stage}
\end{equation}
Here, $\mathbf{X}_t$ is the noisy motion, $c$ collects the text and motion conditions, and $\widehat{\mathbf{R}}_0$ and $\widehat{\mathbf{B}}_0$ are the predicted clean root and body motions. The denoiser first predicts the global root motion, then uses its local representation to condition body prediction. Text embeddings provide semantic guidance, while masked motion constraints support control through keyframes, joint positions and orientations, and root waypoints or paths. In this work, we primarily use the Kimodo-G1 motion generation model as our pretrained backbone, which is trained on human motion-capture data retargeted to the G1 skeleton using the SOMA retargeter to enable direct generation of kinematic robot motions.

\textbf{Motion Representation.} Kimodo~\cite{kimodo} encodes a motion sequence as $\mathbf{X}=[\mathbf{x}_1,\ldots,\mathbf{x}_N]$ for $N$ frames, with the root at the first frame centered horizontally at the origin. For a skeleton with $J$ joints, each frame contains six groups of kinematic features:
\begin{equation}
  \mathbf{x}_n = [\mathbf{r}^{p}_n,\mathbf{r}^{a}_n,\mathbf{j}^{p}_n,\mathbf{j}^{v}_n,\mathbf{j}^{a}_n,\mathbf{f}_n]
  \in \mathbb{R}^{12J+9}.
  \label{eq:kimodo_motion_representation}
\end{equation}
Here, $\mathbf{r}^{p}_n\in\mathbb{R}^{3}$ is the global root position obtained by smoothing the horizontal $(x,z)$ coordinates of the pelvis while preserving its vertical height, and $\mathbf{r}^{a}_n=[\cos\psi_n,\sin\psi_n]\in\mathbb{R}^{2}$ encodes the global heading angle $\psi_n$. Joint positions $\mathbf{j}^{p}_n\in\mathbb{R}^{3J}$ use horizontal offsets from the smoothed root and absolute vertical coordinates, with axes aligned to the world frame. Joint velocities $\mathbf{j}^{v}_n\in\mathbb{R}^{3J}$ are computed from global joint positions, and joint rotations $\mathbf{j}^{a}_n\in\mathbb{R}^{6J}$ encode each joint's global orientation using a continuous 6D representation. The binary vector $\mathbf{f}_n\in\{0,1\}^{4}$ records contact states for the left heel, left toe, right heel, and right toe. These features form a root component $\mathbf{r}^{\mathrm{glob}}_n=[\mathbf{r}^{p}_n,\mathbf{r}^{a}_n]$ and a body component $\mathbf{b}_n=[\mathbf{j}^{p}_n,\mathbf{j}^{v}_n,\mathbf{j}^{a}_n,\mathbf{f}_n]$, providing an explicit decomposition for root and body motion prediction while supporting spatial constraints on root trajectories, joint positions, and joint orientations.

\subsection{ControlNet}
\label{sec:prelim_controlnet}

\textbf{Architecture Overview.} ControlNet~\cite{zhang2023adding} adapts a pretrained diffusion model to new spatial or semantic conditions while preserving the capabilities acquired by its original backbone. Its central design freezes the pretrained denoiser and adds a trainable control branch with a structurally aligned sequence of layers. The control branch processes the additional condition and produces intermediate features that are injected into the corresponding layers of the frozen branch through zero-initialized projections. For a frozen layer $F_{\ell}$, a trainable control layer $G_{\phi,\ell}$, and condition representation $\mathbf{c}$, the interaction can be written as
\begin{equation}
  \mathbf{h}^{\mathrm{ctrl}}_{\ell+1}
  =G_{\phi,\ell}(\mathbf{h}^{\mathrm{ctrl}}_{\ell},\mathbf{c}),
  \qquad
  \mathbf{h}^{\mathrm{lock}}_{\ell+1}
  =F_{\ell}(\mathbf{h}^{\mathrm{lock}}_{\ell})
  +\mathcal{Z}_{\ell}\!\left(\mathbf{h}^{\mathrm{ctrl}}_{\ell+1}\right),
  \label{eq:controlnet_prelim}
\end{equation}
where $\mathcal{Z}_{\ell}$ is a learnable $1\!\times\!1$ projection initialized to zero. At initialization, the residual path therefore contributes no signal and the adapted model exactly follows the pretrained denoiser. During optimization, gradients update only the control branch and zero projections, allowing the model to learn condition-specific behavior while retaining the original generative prior. This structure is particularly useful when the target dataset is limited or when preserving the motion regularities of a pretrained generator is important. In Kimodo-Policy, we adopt this principle for a transformer-based motion denoiser rather than a convolutional image U-Net~\cite{ronneberger2015unetconvolutionalnetworksbiomedical}. The frozen Kimodo root and body denoisers provide the motion prior, while trainable ControlNet branches receive visual features and inject residual hints at selected denoising layers. The specific visual encoder, dual-path fusion mechanism, and root/body injection strategy are described in Sec.~\ref{sec:method_architecture}.
